\chapter{Symplectic Geometry}\label{symplectic-geometry}
\chapterstatus{complete}

\section{Introduction}\label{symplectic-geometry:section-introduction}

A Riemannian metric is a symmetric nondegenerate bilinear form on each tangent space; a \emph{symplectic form} is an \emph{antisymmetric} one, which moreover is closed as a differential form. The
change of symmetry changes everything. A Riemannian metric has local invariants, the curvature, while by Darboux's theorem all symplectic forms of the same dimension look alike near a point; on the
other hand symplectic forms are rigid globally, and live only on manifolds with nonzero even-dimensional cohomology. Symplectic geometry is the geometry of classical mechanics: the phase space of a
mechanical system is the cotangent bundle $T^*Q$ of its configuration space, which carries a canonical symplectic form, and the evolution of the system is the flow of a Hamiltonian vector field. It is also
the geometry behind Kähler manifolds (Chapter~\ref{kahler}), whose K\"ahler form is a symplectic form compatible with the complex structure, and behind the Hodge-theoretic notion of polarisation.

The chapter goes as follows.
\begin{enumerate}
\item Symplectic linear algebra: $\omega^n \neq 0$, isotropic, coisotropic and Lagrangian subspaces, Lagrangian complements and compatible complex structures
(Section~\ref{symplectic-geometry:section-linear}).
\item Symplectic manifolds: the Liouville volume form, the cohomological obstruction $[\omega]^n \neq 0$, examples (Section~\ref{symplectic-geometry:section-manifolds}).
\item The cotangent bundle: the tautological $1$-form, the canonical symplectic form, Lagrangian graphs of closed $1$-forms, conormal bundles, and the splitting of $T(T^*M)$ into two Lagrangian
subbundles given by a torsion-free connection (Section~\ref{symplectic-geometry:section-cotangent}).
\item Hamiltonian vector fields, Poisson brackets and Liouville's theorem; area on the sphere and Archimedes' theorem (Section~\ref{symplectic-geometry:section-hamiltonian}).
\item Darboux's theorem and Moser's theorem on area forms, by Moser's method (Section~\ref{symplectic-geometry:section-darboux}).
\item Lagrangian submanifolds and Weinstein's neighbourhood theorem (Section~\ref{symplectic-geometry:section-lagrangian}).
\end{enumerate}
References: \cite{cannas-da-silva}, \cite{mcduff-salamon}, \cite[Chapter 22]{lee-smooth-manifolds}.

\noindent\textbf{Prerequisites.} Chapter~\ref{linear-algebra} (Linear Algebra), Chapter~\ref{differential-forms} (Differential Forms and Integration), Chapter~\ref{de-rham} (de Rham Cohomology),
Chapter~\ref{connections} (Connections and Curvature) and Chapter~\ref{riemannian-geometry} (Riemannian Geometry).

\section{Symplectic linear algebra}\label{symplectic-geometry:section-linear}

\begin{definition}\label{symplectic-geometry:definition-symplectic-vector-space}
A \emph{symplectic vector space} is a finite-dimensional real vector space $V$ with a nondegenerate alternating bilinear form $\omega$. For a subspace $W$, its \emph{symplectic complement} is
$W^\omega = \{v \in V : \omega(v, w) = 0 \text{ for all } w \in W\}$. The subspace $W$ is \emph{isotropic} if $W \subseteq W^\omega$ (that is, $\omega|_W = 0$), \emph{coisotropic} if $W^\omega \subseteq W$,
\emph{Lagrangian} if $W = W^\omega$, and \emph{symplectic} if $W \cap W^\omega = 0$. A linear map preserving $\omega$ is \emph{symplectic}, and the symplectic automorphisms of $V$ form the
\emph{symplectic group} $\mathrm{Sp}(V, \omega)$.
\end{definition}

By Theorem~\ref{linear-algebra:theorem-symplectic-basis}, $\dim V = 2n$ is even and $V$ has a symplectic basis $e_1, \ldots, e_n, f_1, \ldots, f_n$; with the dual basis,
$\omega = \sum_ie_i^* \wedge f_i^*$. The model is $\RR^{2n}$ with coordinates $(x, y)$ and $\omega_0 = \sum_idx^i \wedge dy^i$, that is, $\omega_0(u, v) = u^TJ_0v$ with
$J_0 = \bigl(\begin{smallmatrix}0 & I\\ -I & 0\end{smallmatrix}\bigr)$; its symplectic group is $\mathrm{Sp}_{2n}(\RR) = \{A : A^TJ_0A = J_0\}$, whose Lie algebra is $\mathfrak{sp}_{2n}(\RR)$
(Example~\ref{representation-theory:example-classical}). In contrast to inner products, there is no notion of length: $\omega(v, v) = 0$ for every $v$.

\begin{proposition}\label{symplectic-geometry:proposition-linear}
Let $(V, \omega)$ be a symplectic vector space of dimension $2n$ and $W \subseteq V$ a subspace.
\begin{enumerate}
\item An alternating form $\sigma$ on a $2n$-dimensional space is nondegenerate if and only if $\sigma^n \neq 0$. Symplectic maps have determinant $1$.
\item $\dim W + \dim W^\omega = 2n$ and $(W^\omega)^\omega = W$.
\item Isotropic subspaces have dimension at most $n$, and a subspace is Lagrangian if and only if it is isotropic of dimension $n$. Every isotropic subspace is contained in a Lagrangian one.
\item Every Lagrangian subspace $L$ has a Lagrangian complement $L^{\prime}$, and then $v \mapsto \omega(\cdot, v)|_L$ is an isomorphism $L^{\prime} \to L^*$; a basis $e_1, \ldots, e_n$ of $L$ and the dual basis
$f_1, \ldots, f_n$ of $L^{\prime}$ ($\omega(e_i, f_j) = \delta_{ij}$) form a symplectic basis. So $(V, \omega) \cong (L \oplus L^*, \omega_L)$ with $\omega_L((u, \alpha), (v, \beta)) = \beta(u) - \alpha(v)$.
\item There is a complex structure $J$ with $\omega(Ju, Jv) = \omega(u, v)$ and $g(u, v) = \omega(u, Jv)$ an inner product (a \emph{compatible} complex structure). For such $J$ and a Lagrangian $L$,
$JL$ is a Lagrangian complement of $L$, equal to the orthogonal complement $L^\perp$ for $g$.
\end{enumerate}
\end{proposition}

\begin{proof}
(1) If $\sigma$ is nondegenerate, $\sigma^n \neq 0$ by Proposition~\ref{multilinear-algebra:proposition-symplectic-determinant}, which also gives the statement on determinants. If $\sigma$ is degenerate,
choose $v \neq 0$ with $\iota_v\sigma = 0$; as $\iota_v$ is an antiderivation, $\iota_v(\sigma^n) = n(\iota_v\sigma) \wedge \sigma^{n-1} = 0$, so the top form $\sigma^n$ vanishes on every basis containing $v$,
hence $\sigma^n = 0$.

(2) The map $V \to W^*$, $v \mapsto \omega(v, \cdot)|_W$, is the composite of the isomorphism $V \to V^*$ given by $\omega$ and the surjective restriction $V^* \to W^*$; its kernel is $W^\omega$. So
$\dim W^\omega = 2n - \dim W$. Clearly $W \subseteq (W^\omega)^\omega$, and both have the same dimension.

(3) If $W \subseteq W^\omega$, then $2\dim W \leq \dim W + \dim W^\omega = 2n$. An isotropic $W$ of dimension $n$ has $\dim W^\omega = n$, so $W = W^\omega$; conversely a Lagrangian is isotropic of dimension
$n$ by (2). If $W$ is isotropic with $\dim W < n$, then $\dim W^\omega > n > \dim W$, so there is $v \in W^\omega \setminus W$, and $W + \RR v$ is isotropic since $\omega(v, v) = 0$; repeat.

(5) Let $e_i$, $f_i$ be a symplectic basis and $Je_i = f_i$, $Jf_i = -e_i$. Then $J^2 = -1$, $J$ preserves $\omega$ (check on basis vectors), and $\omega(u, Jv)$ is the standard inner product in this basis:
$\omega(e_i, Je_j) = \omega(e_i, f_j) = \delta_{ij}$, $\omega(f_i, Jf_j) = -\omega(f_i, e_j) = \delta_{ij}$, and the mixed terms vanish. Given compatible $J$ and Lagrangian $L$, $JL$ is isotropic since
$\omega(Ju, Jv) = \omega(u, v) = 0$, hence Lagrangian; and for $u, v \in L$, $g(u, Jv) = \omega(u, J^2v) = -\omega(u, v) = 0$, so $JL \subseteq L^\perp$, with equality by dimension. In particular
$L \cap JL = 0$.

(4) A Lagrangian complement exists by (5). For $v \in L^{\prime}$, if $\omega(u, v) = 0$ for all $u \in L$, then $v \in L^\omega \cap L^{\prime} = L \cap L^{\prime} = 0$; so $L^{\prime} \to L^*$ is injective, hence an isomorphism.
The dual basis $f_j$ of $L^{\prime}$ satisfies $\omega(e_i, f_j) = \delta_{ij}$, and $\omega(e_i, e_j) = 0 = \omega(f_i, f_j)$ as $L$ and $L^{\prime}$ are isotropic. Identifying $L^{\prime}$ with $L^*$ by
$f \mapsto \omega(\cdot, f)|_L$, for $u, v \in L$ and $f, f^{\prime} \in L^{\prime}$ corresponding to $\alpha$, $\beta$: $\omega(u + f, v + f^{\prime}) = \omega(u, f^{\prime}) - \omega(v, f) = \beta(u) - \alpha(v)$.
\end{proof}

\begin{example}\label{symplectic-geometry:example-linear}
\begin{enumerate}
\item In $(\RR^2, dx \wedge dy)$ every line is Lagrangian, and $\omega(u, v) = \det(u, v)$ is the signed area of the parallelogram spanned by $u$ and $v$. Symplectic maps of $\RR^2$ are the area- and
orientation-preserving linear maps: $\mathrm{Sp}_2(\RR) = \mathrm{SL}_2(\RR)$. In higher dimensions $\mathrm{Sp}_{2n}(\RR) \subsetneq \mathrm{SL}_{2n}(\RR)$: preserving $\omega$ is much stronger than
preserving the volume $\omega^n/n!$.
\item In $(\RR^4, dx^1 \wedge dy^1 + dx^2 \wedge dy^2)$ the planes $\{y = 0\}$, $\{x = 0\}$ and the graph $\{y = Sx\}$ of a symmetric matrix $S$ are Lagrangian; the graph of $A \colon \RR^2 \to \RR^2$ is
Lagrangian if and only if $A$ is symmetric, since $\omega((x, Ax), (x^{\prime}, Ax^{\prime})) = x^TAx^{\prime} - x^{\prime T}Ax$. The plane $\{x^2 = y^2 = 0\}$ is symplectic, and the hyperplane
$\{y^2 = 0\}$ is coisotropic, with $\{y^2 = 0\}^\omega = \RR\partial_{x^2}$.
\item On $\CC^n = \RR^{2n}$ with $z = x + iy$, $\omega_0 = \sum dx^i \wedge dy^i = \frac i2\sum dz^i \wedge d\bar z^i$, and $J$ = multiplication by $i$ is compatible: $\omega_0(u, iv) = \operatorname{Re}\langle u, v\rangle$
for the standard Hermitian product (Proposition~\ref{linear-algebra:proposition-compatible-metric}). The real subspace $\RR^n$ is Lagrangian, and $i\RR^n$ is its orthogonal Lagrangian complement.
\end{enumerate}
\end{example}

\section{Symplectic manifolds}\label{symplectic-geometry:section-manifolds}

\begin{definition}\label{symplectic-geometry:definition-symplectic-manifold}
A \emph{symplectic form} on a manifold $M$ is a closed $2$-form $\omega$ that is nondegenerate at every point; $(M, \omega)$ is a \emph{symplectic manifold}. A \emph{symplectomorphism} is a
diffeomorphism $\varphi$ with $\varphi^*\omega = \omega$. A submanifold $S \subseteq M$ is isotropic, coisotropic, Lagrangian or symplectic if each $T_pS \subseteq T_pM$ is.
\end{definition}

Each tangent space is a symplectic vector space, so $\dim M = 2n$ is even. By Proposition~\ref{symplectic-geometry:proposition-linear}(1), $\omega^n$ vanishes nowhere: \emph{$\omega^n/n!$ is a volume
form}, the \emph{Liouville form}, and every symplectic manifold is oriented by it. Nondegeneracy means that $X \mapsto \iota_X\omega$ is an isomorphism of vector bundles $TM \to T^*M$, an analogue of the
musical isomorphism of a metric (Remark~\ref{riemannian-geometry:remark-musical}), which identifies vector fields with $1$-forms.

\begin{proposition}\label{symplectic-geometry:proposition-cohomology}
If $(M, \omega)$ is a compact symplectic manifold without boundary of dimension $2n$, then $[\omega]^k \neq 0$ in $H^{2k}_{\mathrm{dR}}(M)$ for $0 \leq k \leq n$. In particular $H^{2k}_{\mathrm{dR}}(M) \neq 0$, and
the spheres $S^{2n}$ with $n \geq 2$ carry no symplectic form.
\end{proposition}

\begin{proof}
With the orientation given by $\omega^n$, $\int_M\omega^n > 0$ (Corollary~\ref{differential-forms:corollary-volume}), so $\omega^n$ is not exact by Stokes' theorem
(Theorem~\ref{differential-forms:theorem-stokes}), and $[\omega]^n = [\omega^n] \neq 0$. Since $[\omega]^n = [\omega]^k[\omega]^{n-k}$, also $[\omega]^k \neq 0$. For $n \geq 2$, $H^2_{\mathrm{dR}}(S^{2n}) = 0$
(Example~\ref{de-rham:example-spheres}).
\end{proof}

\begin{example}\label{symplectic-geometry:example-manifolds}
\begin{enumerate}
\item $(\RR^{2n}, \omega_0)$, and every open subset of it.
\item \emph{Surfaces.} On an oriented surface, an area form is a symplectic form: it is closed for degree reasons, and a $2$-form on a $2$-dimensional space is nondegenerate if and only if it is nonzero.
So every orientable surface is symplectic, and symplectic forms on surfaces are just area forms. Symplectomorphisms of surfaces are the area-preserving diffeomorphisms.
\item \emph{Products.} If $(M_1, \omega_1)$, $(M_2, \omega_2)$ are symplectic, so is $M_1 \times M_2$ with $\mathrm{pr}_1^*\omega_1 + \mathrm{pr}_2^*\omega_2$; for instance the torus $T^{2n}$ with the form
induced from $\omega_0$, and $S^2 \times S^2$. On the other hand $S^2 \times S^4$ is not symplectic, although it is orientable with $H^2 \neq 0$: by the K\"unneth formula
(Theorem~\ref{de-rham:theorem-kunneth}) a symplectic form $\omega$ would have $[\omega] = c\,a$ with $a$ the pullback of a generator of $H^2(S^2)$, and $a^2 = 0$, as it is pulled back from
$H^4(S^2) = 0$; so $[\omega]^3 = 0$, contradicting Proposition~\ref{symplectic-geometry:proposition-cohomology}.
\item \emph{K\"ahler manifolds}, such as $\CC\PP^n$ with the Fubini--Study form and all complex projective manifolds, are symplectic (Chapter~\ref{kahler}).
\item \emph{Cotangent bundles} carry a canonical symplectic form (Section~\ref{symplectic-geometry:section-cotangent}).
\end{enumerate}
\end{example}

\section{The cotangent bundle}\label{symplectic-geometry:section-cotangent}

Let $M$ be a manifold of dimension $n$ and $\pi \colon T^*M \to M$ its cotangent bundle. A chart with coordinates $q^1, \ldots, q^n$ on $U \subseteq M$ gives coordinates $(q, p)$ on $T^*U$, where a covector
$\xi = \sum_ip_i\,dq^i|_q$ has coordinates $(q^1, \ldots, q^n, p_1, \ldots, p_n)$.

\begin{definition}\label{symplectic-geometry:definition-tautological}
The \emph{tautological} (or \emph{Liouville}) $1$-form $\lambda$ on $T^*M$ is defined at a point $\xi \in T^*_qM$ by
\[
\lambda_\xi(Y) = \xi\bigl(d\pi(Y)\bigr), \qquad Y \in T_\xi(T^*M);
\]
a point of $T^*M$ is a covector on $M$, and $\lambda$ applies it to the projection of $Y$ to $M$. The \emph{canonical symplectic form} is $\omega = -d\lambda$.
\end{definition}

\begin{proposition}\label{symplectic-geometry:proposition-tautological}
\begin{enumerate}
\item In coordinates $\lambda = \sum_ip_i\,dq^i$ and $\omega = \sum_idq^i \wedge dp_i$; so $\omega$ is a symplectic form.
\item $\lambda$ is characterised by the property $\alpha^*\lambda = \alpha$ for every $1$-form $\alpha$ on $M$, regarded as a section $\alpha \colon M \to T^*M$. Consequently $\alpha^*\omega = -d\alpha$.
\item For a diffeomorphism $f \colon M \to N$, the \emph{cotangent lift} $F \colon T^*M \to T^*N$, $F(\xi) = \xi \circ (d_qf)^{-1}$ for $\xi \in T^*_qM$, satisfies $F^*\lambda_N = \lambda_M$; so it is a
symplectomorphism.
\end{enumerate}
\end{proposition}

\begin{proof}
(1) A tangent vector at $\xi = \sum p_i\,dq^i$ is $Y = \sum_i(a^i\partial_{q^i} + b_i\partial_{p_i})$, and $d\pi(Y) = \sum_ia^i\partial_{q^i}$, so $\lambda(Y) = \sum_ip_ia^i$: $\lambda = \sum_ip_i\,dq^i$. Then
$\omega = -\sum_idp_i \wedge dq^i = \sum_idq^i \wedge dp_i$, which is closed and, at each point, the standard form in the basis $\partial_{q^i}$, $\partial_{p_i}$.

(2) Since $\pi \circ \alpha = \id$, $(\alpha^*\lambda)_q(v) = \lambda_{\alpha_q}(d\alpha(v)) = \alpha_q(d\pi\,d\alpha(v)) = \alpha_q(v)$. Conversely, let $\mu = \sum_i(A_i\,dq^i + B^i\,dp_i)$ in coordinates
satisfy $\alpha^*\mu = \alpha$ for all $\alpha = \sum_ia_i\,dq^i$; then $\sum_iA_i(q, a(q))\,dq^i + \sum_kB^k(q, a(q))\,da_k = \sum_ia_i\,dq^i$. Constant $a_i = c_i$ give $A_i(q, c) = c_i$; then at a point $q_0$,
the functions $a_k = c_k + \sum_jd_{kj}(q^j - q_0^j)$ give $\sum_kB^k(q_0, c)d_{ki} = 0$ for every matrix $(d_{kj})$, so $B^k = 0$. Hence $\mu = \lambda$. Finally $\alpha^*\omega = -d\alpha^*\lambda = -d\alpha$.

(3) $\pi_N \circ F = f \circ \pi_M$, so $(F^*\lambda_N)_\xi(Y) = F(\xi)(d\pi_N\,dF(Y)) = F(\xi)(df\,d\pi_M(Y)) = \xi(d\pi_M(Y)) = \lambda_M(Y)$.
\end{proof}

\begin{corollary}\label{symplectic-geometry:corollary-lagrangian-graphs}
\begin{enumerate}
\item The graph $\alpha(M)$ of a $1$-form $\alpha$ is a Lagrangian submanifold of $T^*M$ if and only if $d\alpha = 0$. In particular the zero section and the graphs $df(M)$ of differentials of functions are
Lagrangian.
\item The fibres $T^*_qM$ are Lagrangian.
\item For a submanifold $S \subseteq M$, the \emph{conormal bundle} $N^*S = \{\xi \in T^*_qM : q \in S, \ \xi|_{T_qS} = 0\}$ is a Lagrangian submanifold of $T^*M$.
\end{enumerate}
\end{corollary}

\begin{proof}
(1) $\alpha \colon M \to T^*M$ is an embedding onto its graph, of dimension $n$, and the restriction of $\omega$ to it corresponds to $\alpha^*\omega = -d\alpha$. (2) In coordinates the fibre is
$\{q = \text{const}\}$, on which every $dq^i$ vanishes. (3) In a slice chart with $S = \{q^{k+1} = \cdots = q^n = 0\}$, $N^*S = \{q^{k+1} = \cdots = q^n = 0, \ p_1 = \cdots = p_k = 0\}$, a submanifold of dimension
$n$ with coordinates $q^1, \ldots, q^k, p_{k+1}, \ldots, p_n$, on which each term $dq^i \wedge dp_i$ vanishes.
\end{proof}

Along the zero section $0_M$ the tangent spaces split canonically: $T_{0_q}(T^*M) = T_q(0_M) \oplus T_{0_q}(T^*_qM) \cong T_qM \oplus T^*_qM$, the second summand being the vertical subspace
(Definition~\ref{connections:definition-vertical}). Both summands are Lagrangian, and in coordinates $\omega(\sum a^i\partial_{q^i}, \sum b_i\partial_{p_i}) = \sum a^ib_i$: the symplectic form pairs them by the
duality $T_qM \times T_q^*M \to \RR$, as in Proposition~\ref{symplectic-geometry:proposition-linear}(4). Away from the zero section the vertical subspace is still canonical, but a horizontal complement requires a
choice, a connection (Section~\ref{connections:section-horizontal}).

\begin{proposition}\label{symplectic-geometry:proposition-connection-lagrangian}
Let $\nabla$ be a connection on $TM$, with the induced connection on $T^*M$ and its horizontal distribution $H \subseteq T(T^*M)$ (Proposition~\ref{connections:proposition-horizontal}). The vertical
distribution $V$ is Lagrangian, and $H$ is Lagrangian if and only if $\nabla$ is torsion-free. In that case $T(T^*M) = H \oplus V$ with $H_\xi \cong T_qM$, $V_\xi \cong T^*_qM$ and
$\omega(h, v) = \langle v, d\pi(h)\rangle$ for $h \in H_\xi$ and $v \in V_\xi$ (identified with a covector by the vertical lift).
\end{proposition}

\begin{proof}
$V$ is the tangent bundle of the fibres, Lagrangian by Corollary~\ref{symplectic-geometry:corollary-lagrangian-graphs}(2). With $\nabla_{\partial_i}\partial_k = \sum_j\Gamma^j_{ik}\partial_j$, the dual connection
satisfies $\nabla_{\partial_i}dq^j = -\sum_k\Gamma^j_{ik}\,dq^k$, so for $\alpha = \sum_jp_j\,dq^j$, $\nabla_{\partial_i}\alpha = \sum_k(\partial_ip_k - \sum_j\Gamma^j_{ik}p_j)\,dq^k$. By
Proposition~\ref{connections:proposition-horizontal}, $H$ is spanned by the horizontal lifts of the coordinate fields,
\[
X_i = \partial_{q^i} + \sum_{j,k}\Gamma^j_{ik}p_j\,\partial_{p_k}.
\]
Then $\omega(X_i, X_l) = dp_i(X_l) - dp_l(X_i) = \sum_jp_j(\Gamma^j_{li} - \Gamma^j_{il}) = -\sum_jp_jT^j_{il}$, where $T(\partial_i, \partial_l) = \sum_jT^j_{il}\partial_j$ is the torsion
(Exercise~\ref{connections:exercise-torsion}). This vanishes for all $p$ if and only if $T = 0$; and $H$ has dimension $n$. Finally $\omega(X_i, \partial_{p_k}) = \delta_{ik}$, which is the stated pairing.
\end{proof}

So a torsion-free connection, for instance the Levi-Civita connection of a metric, splits $T(T^*M)$ into two transverse Lagrangian subbundles paired by duality, extending the canonical splitting along the
zero section, where every horizontal distribution is tangent to the zero section.

\section{Hamiltonian vector fields}\label{symplectic-geometry:section-hamiltonian}

\begin{definition}\label{symplectic-geometry:definition-hamiltonian}
Let $(M, \omega)$ be symplectic. The \emph{Hamiltonian vector field} of a function $H$ is the unique vector field $X_H$ with $\iota_{X_H}\omega = dH$. The \emph{Poisson bracket} of two functions is
$\{F, G\} = \omega(X_F, X_G) = dF(X_G) = X_GF$.
\end{definition}

On $T^*\RR^n$ with $\omega = \sum dq^i \wedge dp_i$, $X_H = \sum_i\bigl(\frac{\partial H}{\partial p_i}\partial_{q^i} - \frac{\partial H}{\partial q^i}\partial_{p_i}\bigr)$, since then
$\iota_{X_H}\omega = \sum_i(\partial_{p_i}H\,dp_i + \partial_{q^i}H\,dq^i) = dH$. So the integral curves of $X_H$ solve \emph{Hamilton's equations}
\[
\dot q^i = \frac{\partial H}{\partial p_i}, \qquad \dot p_i = -\frac{\partial H}{\partial q^i},
\]
and $\{F, G\} = \sum_i(\partial_{q^i}F\,\partial_{p_i}G - \partial_{p_i}F\,\partial_{q^i}G)$.

\begin{lemma}\label{symplectic-geometry:lemma-interior-bracket}
For vector fields $X$, $Y$ on a manifold, $\iota_{[X, Y]} = L_X\iota_Y - \iota_YL_X$ on differential forms.
\end{lemma}

\begin{proof}
Both sides are local antiderivations of degree $-1$: the right side is the commutator of the derivation $L_X$ of degree $0$ (Proposition~\ref{differential-forms:proposition-cartan}) with the antiderivation
$\iota_Y$, and the antiderivation rules combine as in Step 2 of that proof. Both vanish on functions, and on an exact $1$-form $dg$ the left side is $[X, Y]g$ and the right side is
$L_X(Yg) - \iota_YdL_Xg = XYg - YXg$. As in Step 3 of the proof of Proposition~\ref{differential-forms:proposition-cartan}, they agree everywhere.
\end{proof}

\begin{proposition}\label{symplectic-geometry:proposition-hamiltonian}
Let $H$, $F$, $G$ be functions on a symplectic manifold.
\begin{enumerate}
\item $H$ is constant along the integral curves of $X_H$ (conservation of energy), and the flow of $X_H$ preserves $\omega$ and the Liouville volume (Liouville's theorem).
\item $\{F, G\} = -\{G, F\}$, the bracket satisfies the Jacobi identity and the Leibniz rule $\{F, GK\} = \{F, G\}K + G\{F, K\}$, and $[X_F, X_G] = -X_{\{F, G\}}$.
\item $F$ is constant along the flow of $X_G$ if and only if $\{F, G\} = 0$, if and only if $G$ is constant along the flow of $X_F$.
\end{enumerate}
\end{proposition}

\begin{proof}
(1) $X_HH = dH(X_H) = \omega(X_H, X_H) = 0$. By Cartan's formula, $L_{X_H}\omega = d\iota_{X_H}\omega + \iota_{X_H}d\omega = d\,dH = 0$, so the flow preserves $\omega$
(Proposition~\ref{differential-forms:proposition-cartan}) and hence $\omega^n$.

(2) Antisymmetry is that of $\omega$, and the Leibniz rule holds since $X_G$ is linear in $G$ and $\{F, G\} = -X_FG$ is a derivation in $G$. By Lemma~\ref{symplectic-geometry:lemma-interior-bracket},
$\iota_{[X_F, X_G]}\omega = L_{X_F}\iota_{X_G}\omega - \iota_{X_G}L_{X_F}\omega = L_{X_F}dG = d(X_FG) = -d\{F, G\}$, so $[X_F, X_G] = -X_{\{F, G\}}$. For the Jacobi identity, use $\{A, B\} = X_BA$:
$\{\{F, G\}, K\} = X_KX_GF$, $\{\{K, F\}, G\} = X_G\{K, F\} = -X_GX_KF$, and $\{\{G, K\}, F\} = -\{F, \{G, K\}\} = -X_{\{G, K\}}F = [X_G, X_K]F$. The sum is $[X_K, X_G]F + [X_G, X_K]F = 0$.

(3) $\frac{d}{dt}F(\varphi_t(p)) = (X_GF)(\varphi_t(p)) = \{F, G\}(\varphi_t(p))$, and antisymmetry.
\end{proof}

\begin{example}\label{symplectic-geometry:example-hamiltonian}
\begin{enumerate}
\item \emph{The harmonic oscillator} $H = \frac12(q^2 + p^2)$ on $\RR^2$: Hamilton's equations $\dot q = p$, $\dot p = -q$ give the rotations $(q, p) \mapsto (q\cos t + p\sin t, -q\sin t + p\cos t)$, which preserve
$H$ (circles) and area.
\item \emph{The pendulum} $H = \frac12p^2 - \cos q$ on $T^*S^1$: $\dot q = p$, $\ddot q = \dot p = -\sin q$. The level sets of $H$ are the phase curves: closed curves around $(0, 0)$ for $H < 1$ (oscillations), curves
going around the cylinder for $H > 1$ (rotations), and the separatrix $H = 1$ through the unstable equilibrium $(\pi, 0)$.
\item \emph{Geodesic flow.} On $T^*M$ for a Riemannian manifold, $H(\xi) = \frac12|\xi|^2$ (the metric on covectors, Remark~\ref{riemannian-geometry:remark-musical}): in coordinates $H = \frac12\sum g^{ij}p_ip_j$, and
Hamilton's equations are $\dot q^k = \sum_ig^{ki}p_i$, $\dot p_k = -\frac12\sum_{i,j}\partial_kg^{ij}p_ip_j$. Eliminating $p$ by $p_i = \sum_jg_{ij}\dot q^j$ and using
$\partial_kg^{ij} = -\sum_{a,b}g^{ia}\partial_kg_{ab}g^{bj}$, they become $\frac{d}{dt}\sum_jg_{kj}\dot q^j = \frac12\sum_{a,b}\partial_kg_{ab}\dot q^a\dot q^b$, which, expanding the left side and multiplying by
$(g^{lk})$, is the geodesic equation $\ddot q^l + \sum\Gamma^l_{ab}\dot q^a\dot q^b = 0$ (Theorem~\ref{riemannian-geometry:theorem-levi-civita}). So
geodesics are the projections of the integral curves of a Hamiltonian vector field, and the geodesic flow preserves the Liouville measure of $T^*M$.
\end{enumerate}
\end{example}

\begin{example}\label{symplectic-geometry:example-archimedes}
\emph{Area on the sphere and Archimedes' theorem.} On $S^2 \subseteq \RR^3$ use the cylindrical coordinates $(\theta, z) \in (0, 2\pi) \times (-1, 1)$, $x = (r\cos\theta, r\sin\theta, z)$ with $r = \sqrt{1 - z^2}$.
Then $x_\theta = (-r\sin\theta, r\cos\theta, 0)$ and $x_z = (-\frac zr\cos\theta, -\frac zr\sin\theta, 1)$ are orthogonal, with lengths $r$ and $1/r$, and $x_\theta \times x_z = x$ is the outer normal. So
the area form of the outward orientation is
\[
dA = |x_\theta \times x_z|\,d\theta \wedge dz = d\theta \wedge dz :
\]
the projection from the sphere to the cylinder $\{r = 1\}$ along horizontal rays from the axis preserves area. This is \emph{Archimedes' theorem}; in particular the zone $\{a < z < b\}$ has area $2\pi(b - a)$,
independently of where it is, and the whole sphere has area $4\pi$. For the symplectic form $dA$, $\iota_{\partial_\theta}(d\theta \wedge dz) = dz$, so the Hamiltonian vector field of the height function $z$
is $\partial_\theta$, the rotation about the vertical axis: the height is the conserved quantity of the rotation, its \emph{moment map}.
\end{example}

\section{Darboux's theorem and Moser's method}\label{symplectic-geometry:section-darboux}

Let $X_t$, $t \in [0, 1]$, be a time-dependent vector field, smooth in $(t, x)$, and $\varphi_t$ its flow, $\frac{d}{dt}\varphi_t(x) = X_t(\varphi_t(x))$ with $\varphi_0 = \id$, defined where the solutions exist
(Theorem~\ref{real-analysis:theorem-picard-lindelof}). For a smooth family of forms $\omega_t$,
\begin{equation}\label{symplectic-geometry:equation-moser}
\frac{d}{dt}\varphi_t^*\omega_t = \varphi_t^*\Bigl(L_{X_t}\omega_t + \frac{d\omega_t}{dt}\Bigr);
\end{equation}
this is the time-dependent version of the formula $\frac{d}{dt}\Phi_t^*\omega = \Phi_t^*L_X\omega$ of Proposition~\ref{differential-forms:proposition-cartan}, with the same proof \cite[Chapter 22]{lee-smooth-manifolds}.
\emph{Moser's method} constructs a diffeomorphism $\varphi_1$ with $\varphi_1^*\omega_1 = \omega_0$ by choosing $X_t$ so that the right side vanishes.

\begin{theorem}[Darboux]\label{symplectic-geometry:theorem-darboux}
Every point of a symplectic manifold $(M, \omega)$ of dimension $2n$ has a chart with coordinates $(x^1, \ldots, x^n, y^1, \ldots, y^n)$ in which $\omega = \sum_idx^i \wedge dy^i$.
\end{theorem}

\begin{proof}[Sketch of proof]
Choose a chart centred at the given point; after a linear change of coordinates (Theorem~\ref{linear-algebra:theorem-symplectic-basis}) the form $\omega_1 = \omega$ agrees at $0$ with the constant form
$\omega_0 = \sum dx^i \wedge dy^i$. Let $\omega_t = \omega_0 + t(\omega_1 - \omega_0)$. At $0$ all $\omega_t$ equal $\omega_0$, so by compactness of $[0, 1]$ and continuity of the determinant they are
nondegenerate on a ball $B$ around $0$ for all $t \in [0, 1]$. The closed form $\omega_1 - \omega_0$ is $d\beta$ on $B$ (Theorem~\ref{de-rham:theorem-homotopy-invariance}), and subtracting a constant
$1$-form we may assume $\beta(0) = 0$. By nondegeneracy there is a unique $X_t$ with $\iota_{X_t}\omega_t = -\beta$; it is smooth in $(t, x)$ and vanishes at $0$, so its flow fixes $0$ and, by continuous
dependence on initial conditions, is defined for $t \in [0, 1]$ on a smaller ball. By \eqref{symplectic-geometry:equation-moser} and Cartan's formula,
\[
\frac{d}{dt}\varphi_t^*\omega_t = \varphi_t^*\bigl(d\iota_{X_t}\omega_t + \iota_{X_t}d\omega_t + \omega_1 - \omega_0\bigr) = \varphi_t^*(-d\beta + 0 + d\beta) = 0 .
\]
So $\varphi_1^*\omega_1 = \omega_0$, and composing the chart with $\varphi_1$ gives the required coordinates. The omitted verification is formula \eqref{symplectic-geometry:equation-moser}.
\end{proof}

So symplectic manifolds have no local invariants: there is no symplectic analogue of curvature. Compare with Riemannian geometry, where a metric is locally Euclidean only if it is flat
(Theorem~\ref{connections:theorem-flat-local-systems} applied to the Levi-Civita connection).

\begin{theorem}[Moser]\label{symplectic-geometry:theorem-moser-area}
Let $\Sigma$ be a compact connected oriented surface and $\omega_0$, $\omega_1$ area forms (positive for the orientation) with $\int_\Sigma\omega_0 = \int_\Sigma\omega_1$. Then there is a diffeomorphism
$\varphi$ of $\Sigma$ with $\varphi^*\omega_1 = \omega_0$. In particular two compact connected orientable surfaces with area forms are symplectomorphic if and only if they are diffeomorphic and have the same area.
\end{theorem}

\begin{proof}[Sketch of proof]
The forms $\omega_t = (1 - t)\omega_0 + t\omega_1$ are positive, as convex combinations of positive top forms, hence area forms. Since $\int_\Sigma(\omega_1 - \omega_0) = 0$, the form $\omega_1 - \omega_0$ is
exact (Corollary~\ref{de-rham:corollary-consequences}), say $d\beta$. Define $X_t$ by $\iota_{X_t}\omega_t = -\beta$; as $\Sigma$ is compact its flow is defined for all $t \in [0, 1]$, and the computation in the
proof of Theorem~\ref{symplectic-geometry:theorem-darboux} gives $\varphi_1^*\omega_1 = \omega_0$. The same argument, with ``area form'' replaced by ``volume form'', applies in every dimension; it again uses
formula \eqref{symplectic-geometry:equation-moser}.
\end{proof}

\section{Lagrangian submanifolds}\label{symplectic-geometry:section-lagrangian}

Lagrangian submanifolds are the most important submanifolds of symplectic manifolds: graphs of closed $1$-forms (Corollary~\ref{symplectic-geometry:corollary-lagrangian-graphs}), and, as the next example
shows, graphs of symplectomorphisms.

\begin{example}\label{symplectic-geometry:example-lagrangian}
\begin{enumerate}
\item Every curve in a surface is Lagrangian: a line in a $2$-dimensional symplectic vector space is Lagrangian.
\item \emph{Graphs of symplectomorphisms.} For a diffeomorphism $\varphi$ of $(M, \omega)$, the graph $\Gamma_\varphi = \{(x, \varphi(x))\} \subseteq M \times M$ is Lagrangian for the form
$\mathrm{pr}_2^*\omega - \mathrm{pr}_1^*\omega$ if and only if $\varphi$ is a symplectomorphism: the parametrisation $x \mapsto (x, \varphi(x))$ pulls the form back to $\varphi^*\omega - \omega$.
\item The torus $\{|z_1| = \cdots = |z_n| = 1\} \subseteq (\CC^n, \omega_0)$ is Lagrangian. Its tangent space at $z$ is spanned by the vectors $iz_je_j$, with $(e_j)$ the standard basis, and since
$\omega_0(u, v) = \operatorname{Re}\langle u, -iv\rangle$ (Example~\ref{symplectic-geometry:example-linear}(3)), $\omega_0(iz_je_j, iz_ke_k) = \operatorname{Re}\langle iz_je_j, z_ke_k\rangle = \delta_{jk}\operatorname{Re}(i|z_j|^2) = 0$.
\end{enumerate}
\end{example}

\begin{theorem}[Weinstein's Lagrangian neighbourhood theorem]\label{symplectic-geometry:theorem-weinstein}
Let $L$ be a compact Lagrangian submanifold of a symplectic manifold $(M, \omega)$. Then a neighbourhood of $L$ in $M$ is symplectomorphic to a neighbourhood of the zero section in $(T^*L, -d\lambda)$, by a
symplectomorphism that is the identity on $L$ (identified with the zero section).
\end{theorem}

This is \cite{cannas-da-silva} (Lagrangian neighbourhood theorem) or \cite{mcduff-salamon}. The proof chooses a compatible almost complex structure, uses $J$ to identify the normal bundle of $L$ with $TL$, and the
metric with $T^*L$ (Proposition~\ref{symplectic-geometry:proposition-linear}(4),(5), pointwise), so that a tubular neighbourhood gives a diffeomorphism which is symplectic along $L$; Moser's method, relative to
$L$, then corrects it to a symplectomorphism. Consequences:
\begin{enumerate}
\item \emph{Nearby Lagrangians.} A Lagrangian submanifold $C^1$-close to $L$ corresponds, under the theorem, to a submanifold of $T^*L$ close to the zero section and transverse to the fibres, hence to the graph of a
$1$-form, which is closed by Corollary~\ref{symplectic-geometry:corollary-lagrangian-graphs}. So the Lagrangians near $L$ are parametrised by the closed $1$-forms on $L$ near $0$. The graphs of exact forms
$df$ are the images of the zero section under the time-one maps of the Hamiltonians $-f \circ \pi$ (Exercise~\ref{symplectic-geometry:exercise-generating-function}, in a chart).
\item \emph{Symplectomorphisms near the identity.} Applied to the diagonal of $M \times M$, the graph of the identity, which is Lagrangian for $\mathrm{pr}_2^*\omega - \mathrm{pr}_1^*\omega$ by
Example~\ref{symplectic-geometry:example-lagrangian}(2), the theorem shows that symplectomorphisms $C^1$-close to the identity correspond to closed $1$-forms on $M$ near $0$; in particular such a
symplectomorphism, corresponding to $\alpha$, is joined to the identity by the path of symplectomorphisms corresponding to $t\alpha$, $0 \leq t \leq 1$.
\end{enumerate}

\section{Exercises}\label{symplectic-geometry:section-exercises}

\begin{exercise}\label{symplectic-geometry:exercise-nondegenerate-r4}
For which real numbers $a$ is $\sigma = dx^1 \wedge dy^1 + a\,dx^1 \wedge dx^2 + dx^2 \wedge dy^2$ a symplectic form on $\RR^4$? Find a Lagrangian plane for it.
\end{exercise}

\begin{solution}
$\sigma$ is closed (constant coefficients). $\sigma^2 = 2\,dx^1 \wedge dy^1 \wedge dx^2 \wedge dy^2$, since $dx^1 \wedge dx^2$ wedges to zero with the other two terms and with itself; so $\sigma$ is symplectic for
every $a$ (Proposition~\ref{symplectic-geometry:proposition-linear}(1)). The plane $\{x^1 = x^2 = 0\}$ is Lagrangian, as every term contains $dx^1$ or $dx^2$.
\end{solution}

\begin{exercise}\label{symplectic-geometry:exercise-archimedes-cap}
Using Example~\ref{symplectic-geometry:example-archimedes}, compute the area of the spherical cap $\{z > \cos\vartheta_0\}$ and compare with the holonomy of the sphere computed in
Example~\ref{connections:example-transport}(2).
\end{exercise}

\begin{solution}
The cap is $\{\cos\vartheta_0 < z < 1\}$, of area $2\pi(1 - \cos\vartheta_0)$; this is the angle of the rotation produced by parallel transport around its boundary circle, modulo $2\pi$, in accordance with the
Gauss--Bonnet theorem for the curvature $1$.
\end{solution}

\begin{exercise}\label{symplectic-geometry:exercise-not-symplectic}
Show that $S^4$, $S^2 \times S^4$ and $S^1 \times S^3$ admit no symplectic form, and that $T^4$, $S^2 \times S^2$ and $S^2 \times T^2$ do.
\end{exercise}

\begin{solution}
$H^2_{\mathrm{dR}}$ vanishes for $S^4$ and for $S^1 \times S^3$ (K\"unneth, Theorem~\ref{de-rham:theorem-kunneth}), and for $S^2 \times S^4$ the square of every class in $H^2$ vanishes, as in
Example~\ref{symplectic-geometry:example-manifolds}(3); Proposition~\ref{symplectic-geometry:proposition-cohomology} applies. Products of symplectic manifolds (tori, area forms on $S^2$) are symplectic.
\end{solution}

\begin{exercise}\label{symplectic-geometry:exercise-generating-function}
Let $f \colon \RR^n \to \RR$ be smooth. Show that $\Phi(q, p) = (q, p + df(q))$ is a symplectomorphism of $T^*\RR^n$, the time-one map of the Hamiltonian $H(q, p) = -f(q)$, which maps the zero section onto the
Lagrangian graph of $df$.
\end{exercise}

\begin{solution}
$\Phi^*\sum dq^i \wedge dp_i = \sum dq^i \wedge (dp_i + \sum_j\partial_i\partial_jf\,dq^j) = \sum dq^i \wedge dp_i$, since $\sum_{i,j}\partial_i\partial_jf\,dq^i \wedge dq^j = 0$ by symmetry. For $H = -f(q)$,
Hamilton's equations are $\dot q = 0$, $\dot p = \nabla f(q)$, with time-one map $\Phi$.
\end{solution}

\begin{exercise}\label{symplectic-geometry:exercise-poisson-r2n}
On $T^*\RR^n$ compute $\{q^i, q^j\}$, $\{p_i, p_j\}$ and $\{q^i, p_j\}$, and show that the angular momentum $L = q^1p_2 - q^2p_1$ Poisson-commutes with every rotation-invariant Hamiltonian
$H = h(|q|^2, |p|^2, \langle q, p\rangle)$ on $T^*\RR^2$.
\end{exercise}

\begin{solution}
With $\{F, G\} = \sum_k(\partial_{q^k}F\,\partial_{p_k}G - \partial_{p_k}F\,\partial_{q^k}G)$: $\{q^i, q^j\} = \{p_i, p_j\} = 0$ and $\{q^i, p_j\} = \delta_{ij}$. The flow of $X_L$ is $\dot q = (-q^2, q^1)$,
$\dot p = (-p_2, p_1)$, the simultaneous rotation of $q$ and $p$, which preserves $|q|^2$, $|p|^2$ and $\langle q, p\rangle$; so $H$ is constant along it, and $\{H, L\} = X_LH = 0$
(Proposition~\ref{symplectic-geometry:proposition-hamiltonian}(3)).
\end{solution}
